\paragraph{抽象\trBraneStructure。}\label{sub:abstract-brane-structures}
数对 $(\mu,w_2)$ 由一个映入\trEilenbergMacLane{}空间乘积的映射所表示%
\footnote{译注：原书此处印作 $K(\ZZ/1)$；疑应为 $K(\ZZ,1)$（$\mu$ 为整值的
Maslov 类，须经 $K(\ZZ,1)$ 表示；且唯有取 $K(\ZZ,1)\times K(\ZZ/2,2)$ 的两个%
\trUniversal\trFibration{}之乘积，\trPullback{}所得\trFiber{}才是下文的
$\ZZ\times\RP^\infty$）；此处径改。
——译者注}，
\begin{equation}\label{eq:maslov-w2-classifying-map}
Gr(V)\longrightarrow K(\ZZ,1)\times K(\ZZ/2,2).
\end{equation}
沿该映射\trPullback{}那两个\trUniversal\trFibration{}的乘积，便得到一个\trFibration{}
$Gr(V)^\#\to Gr(V)$，其\trFiber{}为 $\ZZ\times\RP^\infty$。当 $n\ge 3$ 时，可以把
$Gr(V)^\#$ 等同于 $Gr(V)$ 在\trPostnikov{}分解
意义下的 $2$-\trConnected{}的“\trCovering”。当 $n\le 2$ 时它仍然\trConnected，
并且分别满足
\begin{equation}\label{eq:postnikov-low-dim-homotopy}
\begin{cases}
\pi_1(Gr(V)^\#)\cong\ZZ/2,\ \ \pi_2(Gr(V)^\#)=0 & \text{当 } n=1;\\
\pi_1(Gr(V)^\#)=0,\ \ \pi_2(Gr(V)^\#)\cong\ZZ & \text{当 } n=2.
\end{cases}
\end{equation}
我们把 $Gr(V)^\#$ 的点称为\textbf{抽象\trLinear\trLagrange\trBrane}
（abstract linear Lagrangian branes），记作 $\Lambda^\#$，其中 $\Lambda$ 是底层的%
\trLagrange{}子空间。
\begin{lemma}\label{thm:abstract-brane-loop-determinant}
取一条\trLoop{} $\rho^\#\in\scr{L}\,Gr(V)^\#$，并设 $D_{D,\rho}$ 是\trDisk{} $D$ 上
关联于底层\trLoop{} $\rho$ 的 $\bar\partial$-算子。那么
$\mathrm{index}(D_{D,\rho})=n$；而且 $\rho^\#$ 决定一个\trIsomorphism
\begin{equation*}
\det(D_{D,\rho})\cong\lambda^{\mathrm{top}}(\rho(0)).
\end{equation*}
\end{lemma}
\begin{proof}
能\trLift{}到 $Gr(V)^\#$ 的\trLoop{}必满足 $\mu(\rho)=0$，因此关于\trIndex{}的陈述是
\cref{thm:det-line-bundle-loop-space}中陈述的一个特例。映射
$\scr{L}\,Gr(V)^\#\to\scr{L}_0\,Gr(V)$ 的诸\trFiber{}\trHomotopyEquivalence{}于
$\ZZ\times\scr{L}\,\RP^\infty$。忘掉第一个因子而只看第二个因子的 $\pi_0$，便得到
一个诱导的\trDoubleCover，从而可以配一个\trRealLineBundle{}
$\delta_D\to\scr{L}_0\,Gr(V)$。由构造，它由
$T(w_2)\in H^1(\scr{L}_0\,Gr(V);\ZZ/2)$ 分类，故由
\cref{thm:det-line-bundle-loop-space}，它\trIsomorphic{}于以
\begin{equation}\label{eq:determinant-fibre-line}
\det(D_{D,\rho})^\vee\otimes\lambda^{\mathrm{top}}(\rho(0)).
\end{equation}
为\trFiber{}的自然\trLineBundle。要在两者之间确定一个\trIsomorphism，只需考察
$\scr{L}_0\,Gr(V)$ 的任意\trConnected{}子集即可，例如由常值\trLoop{}组成的那一个。回忆
\cref{thm:det-line-bundle-loop-space}的证明：对常值的 $\rho$，$D_{D,\rho}$ 是%
\trSurjective{}的，其\trKernel{}\trIsomorphic{}于 $\rho(0)$；这意味着
$\det(D_{D,\rho})^\vee\otimes\lambda^{\mathrm{top}}(\rho(0))$ 有一个自然的取定%
\trOrientation。另一方面，$\scr{L}\,Gr(V)^\#$ 中位于 $\rho$ 之上的\trFiber{}可以
等同于 $\ZZ\times\scr{L}\,\RP^\infty$，且此等同在相差\trHomotopy{}与 $\ZZ$ 因子中
的\trTranslation{}的意义下唯一；于是取
$\pi_0(\scr{L}\,\RP^\infty)\cong\pi_1(\RP^\infty)$ 的平凡元素便得到 $(\delta_D)_\rho$
的一个\trOrientation。把二者结合起来便得到所要的\trLineBundle{}\trIsomorphism。最后，由
定义，$\delta_D$ 到 $\scr{L}\,Gr(V)^\#$ 上的\trPullback{}是\trCanonical{}平凡的，
从而我们得到\eqref{eq:determinant-fibre-line} 之\trPullback{}的一个%
\trTrivialization。
\end{proof}
对抽象\trLinear\trLagrange\trBrane{}的任一对
$(\Lambda_0^\#,\Lambda_1^\#)$——只要底层的 $(\Lambda_0,\Lambda_1)$
\trTransverse{}相交——都可以按下述方式关联一个整数，即绝对\trIndex{}
$i(\Lambda_0^\#,\Lambda_1^\#)$，以及一个一维实\trVectorSpace，即\trOrientation{}
空间 $o(\Lambda_0^\#,\Lambda_1^\#)$。在 $Gr(V)^\#$ 中取一条连接这两点的路径
$\lambda^\#$，其\trProjection{} $\lambda$ 落在 $\scr{P}^-\,Gr(V)$ 中，并令
\begin{equation}\label{eq:absolute-index-orientation-space}
i(\Lambda_0^\#,\Lambda_1^\#)=I_H(\lambda),\qquad
o(\Lambda_0^\#,\Lambda_1^\#)=(\delta_H)_\lambda.
\end{equation}
当然，必须证明这些不依赖于 $\lambda^\#$ 的选取（就 $o$ 而言是相差%
\trCanonical\trIsomorphism）。把 $\Omega^-(Gr(V),\Lambda_0,\Lambda_1)$ 记作
$\scr{P}^-\,Gr(V)$ 中由从 $\Lambda_0$ 到 $\Lambda_1$ 的路径组成的子空间，而把
$\Omega^-(Gr(V)^\#,\Lambda_0^\#,\Lambda_1^\#)$ 记作 $Gr(V)^\#$ 中相应的路径
$\lambda^\#$ 的空间。这两个空间分别\trHomotopyEquivalence{}于通常的\trLoop{}空间
$\Omega\,Gr(V)$ 与 $\Omega\,Gr(V)^\#$。当 $n\ge 3$ 时，我们知道
$\Omega\,Gr(V)^\#$ 是\trSimplyConnected{}的，这立即蕴含所要的唯一性结果。当
$n=2$ 时，$w_2$ 在 $\pi_2(Gr(V))$ 上仍然非平凡，这意味着映射
$\ZZ=\pi_2(Gr(V)^\#)\to\pi_2(Gr(V))=\ZZ$ 是乘以 $2$。于是，\trLineBundle{}
$\delta_H$ \trPullback{}到 $\Omega^-(Gr(V)^\#,\Lambda_0^\#,\Lambda_1^\#)$ 上是平凡的，这
同样给出所要的结论。在最后的情形 $n=1$，由
\eqref{eq:postnikov-low-dim-homotopy} 有两个不同的\trHomotopy{}类；二者给出同一个
$i$；至于\trOrientation{}空间，我们索性规定：$\lambda^\#$ 的\trHomotopy{}类的更换
会给出
$(\delta_H)_\lambda$ 的一个\trOrientationReversing{}的\trAutomorphism（一个%
\trEquivalence{}的做法
——也许更系统一些——是用\trStabilization{}来处理低维情形）。
设 $S$ 是一个带\trStriplikeEnd{}的\trRiemann\trSurface，其上配备平凡\trSymplectic
\trVectorBundle{} $E=S\times V$，以及一个由映射 $\nu\colon\partial S\to Gr(V)$
给出的\trLagrange\trSubbundle{} $F_\nu\subset E|_{\partial S}$，该映射在诸%
\trStriplikeEnd{}上是\trLocallyConstant{}的。后一条件意味着：对无穷远的任一点
$\zeta$，都有一对\trLagrange{}子空间 $(\Lambda_{\zeta,0},\Lambda_{\zeta,1})$ 使得
$\nu(\epsilon_\zeta(s,k))=\Lambda_{\zeta,k}$，其中
$(s,k)\in\RR^{\pm}\times\{0;1\}$。我们将假设每一对都满足
$\Lambda_{\zeta,0}\pitchfork\Lambda_{\zeta,1}$，从而关联的 $\bar\partial$-算子
——记作 $D_{S,\nu}$——是\trNondegenerate{}的。此外，设还给定了 $\nu$ 到
$Gr(V)^\#$ 中的一个\trLift{} $\nu^\#$，它在诸端上同样\trLocallyConstant，从而为
诸\trLagrange{}子空间给出特异的\trPreimage{} $\Lambda_{\zeta,k}^\#$。
\begin{proposition}\label{thm:general-surface-index-formula}
把\trCompactification{} $\hat{S}$ 的诸\trBoundaryComponent{}记作
$\hat{C}_1,\ldots,\hat{C}_k$，并令 $e_j=|\Sigma^-\cap\hat{C}_j|$ 为落在每条分支上
的无穷远\trIncomingPoint{}的个数。对每个 $j$ 取定某个非无穷远点
$z_j\in\hat{C}_j$。那么
\begin{equation}\label{eq:general-surface-index-formula}
\begin{gathered}
\mathrm{index}(D_{S,\nu})=n(\chi(S)-|\Sigma^-|)
+\sum_{\zeta\in\Sigma^{\pm}}\mp\,i(\Lambda_{\zeta,0}^\#,\Lambda_{\zeta,1}^\#),\\
\det(D_{S,\nu})\cong\bigotimes_j\Bigl(\lambda^{\mathrm{top}}(\nu_{z_j})^{\otimes(1-e_j)}
\otimes\bigotimes_{\zeta\in\Sigma^+\cap\hat{C}_j}
o(\Lambda_{\zeta,0}^\#,\Lambda_{\zeta,1}^\#)^{\mp 1}\Bigr).
\end{gathered}
\end{equation}
\end{proposition}
\begin{proof}
我们已经证明过两个特例。当 $S=D$ 时，陈述恰好就是
\cref{thm:abstract-brane-loop-determinant}。在另一种情形 $S=H$，首先需要对
$\nu\colon\partial H=\RR\to Gr(V)$ 作一个小的\trCompactlySupported
\trDeformation，把它化成 $\nu(s)=\lambda(\psi(s))$ 的形式，其中
$\lambda\in\Omega^-(Gr(V),\Lambda_{\xi,0},\Lambda_{\xi,1})$；对诸\trLift{}
$\nu^\#$、$\lambda^\#$ 亦然。这诱导了\trFredholmOperator{}的一个\trDeformation，
因而不改变\eqref{eq:general-surface-index-formula} 的左边；而右边显然也保持
不变，因为它不依赖于路径的任何选取。做完这一点后，便可诉诸
\cref{thm:halfplane-index-determinant}以及 $i$、$o$ 的定义。我们还需要一个相关的
观察。取带一个出向无穷远点的\trUpperHalfPlane，分别记作 $\bar{H}$ 与
$\bar{\zeta}$。给它配备相反的路径 $\bar{\nu}^\#(s)=\nu(-s)$，使得关联的%
\trLagrange{}子空间对与从前一致：$\Lambda_{\bar{\xi},k}=\Lambda_{\xi,k}$。于是
可以把 $\bar{H}\#H\cong D$ \trGluing{}起来，这与全部附加结构\trCompatible，故
施用\eqref{eq:delbar-gluing} 便得
\begin{equation}\label{eq:opposite-halfplane-index}
\begin{gathered}
n=\mathrm{index}(D_{\bar{H},\bar{\nu}})
+i(\Lambda_{\xi,0}^\#,\Lambda_{\xi,1}^\#),\\
\lambda^{\mathrm{top}}(\bar{\nu}_z)\cong\det(D_{\bar{H},\bar{\nu}})
\otimes o(\Lambda_{\xi,0}^\#,\Lambda_{\xi,1}^\#)
\end{gathered}
\end{equation}
其中 $z$ 是 $\partial\bar{H}$ 的任一点。这证明了
\eqref{eq:general-surface-index-formula} 对 $\bar{H}$ 也成立。
一般的策略如下。可以把\trCompactification{} $\hat{S}$ 想成 $k$ 个\trDisk{}
$\hat{D}_1,\ldots,\hat{D}_k$ 与一个同\trGenus{} $g$ 的闭\trSurface{}的%
\trConnectedSum。
把这些\trDisk{}的有限（\trPointedBoundary）部分记作 $D_1,\ldots,D_k$。一如
\cref{thm:det-line-bundle-loop-space}证明的第 (i) 步，可以\trPinch{}掉与
$\partial\hat{D}_j$ 平行的曲线，从而证明
\begin{equation*}
\begin{gathered}
\mathrm{index}(D_{S,\nu})=n(2-2g)
+\sum_j\bigl(\mathrm{index}(D_{D_j,\nu_j})-2n\bigr),\\
\det(D_{S,\nu})\cong\bigotimes_j\det(D_{D_j,\nu_j}),
\end{gathered}
\end{equation*}
其中 $\nu_j$ 是 $\nu$ 在属于 $\hat{C}_j$ 的那些分支上的限制。这把问题化归到
$\hat{S}$ 是一个\trDisk{}的情形，故我们从此假设情况正是如此。对每个出向无穷远点
$\xi^+\in\Sigma^+$，在
$\Omega^-(Gr(V)^\#,\Lambda_{\xi^+,0}^\#,\Lambda_{\xi^+,1}^\#)$ 中取一条路径。取
若干配备相应边界数据的 $H$ 的拷贝，把它们\trGluing{}到 $S$ 上以填满诸出端。如果
只有一个无穷远\trIncomingPoint，那么这一\trGluing{}过程的结果是 $H$ 上的另一个
算子，而作为前面已证情形的推论，可以看出%
\footnote{译注：下式第二行左端 $o$ 的第一对变量原书印作
$\Lambda_{\xi^-,0},\Lambda_{\xi^-,1}$，无 $\#$ 上标；$o$ 仅定义于抽象\trLinear\trLagrange
\trBrane{}之对\eqref{eq:absolute-index-orientation-space}，疑应为
$\Lambda_{\xi^-,0}^\#,\Lambda_{\xi^-,1}^\#$；此处照原书排印，存疑备案。
——译者注}
\begin{equation*}
\begin{gathered}
i(\Lambda_{\xi^-,0}^\#,\Lambda_{\xi^-,1}^\#)=\mathrm{index}(D_{S,\nu})
+\sum_{\xi^+}i(\Lambda_{\xi^+,0}^\#,\Lambda_{\xi^+,1}^\#),\\
o(\Lambda_{\xi^-,0},\Lambda_{\xi^-,1})\cong\det(D_{S,\nu})
\otimes\bigotimes_{\xi^+}o(\Lambda_{\xi^+,0}^\#,\Lambda_{\xi^+,1}^\#).
\end{gathered}
\end{equation*}
如果没有、或者有 $\ge 2$ 个无穷远\trIncomingPoint，则按类似却稍复杂的方式进行，
另外再\trGluing{}上若干 $\bar{H}$ 的拷贝把 $S$ 完全\trCompactification，然后利用
\eqref{eq:opposite-halfplane-index} 得出
\begin{equation}\label{eq:compactified-conclusion}
\begin{gathered}
n=\sum_{\xi^-}\bigl(n-i(\Lambda_{\xi^-,0}^\#,\Lambda_{\xi^-,1}^\#)\bigr)
+\mathrm{index}(D_{S,\nu})
+\sum_{\xi^+}i(\Lambda_{\xi^+,0}^\#,\Lambda_{\xi^+,1}^\#),\\
\lambda^{\mathrm{top}}(\nu_{z_0})\cong\bigotimes_{\xi^-}
\bigl(\lambda^{\mathrm{top}}(\nu_{z_j})\otimes
o(\Lambda_{\xi^-,0}^\#,\Lambda_{\xi^-,1}^\#)^{-1}\bigr)
\otimes\det(D_{S,\nu})\\
\otimes\bigotimes_{\xi^+}o(\Lambda_{\xi^+,0}^\#,\Lambda_{\xi^+,1}^\#).
\end{gathered}
\end{equation}
诸 $z_j$ 是 $S$ 的边界点。可以沿\trCompactification\trSurface{}上\trLagrange{}子空间的%
\trLoop{}作\trParallelTransport{}，把诸 $\lambda^{\mathrm{top}}(\nu_{z_j})$ 彼此等同（由于该%
\trLoop{}容许一个\textbf{\trBraneStructure}，其\trMaslovIndex{}为零，从而%
\trOrientable，故这里没有歧义）；这就把\eqref{eq:compactified-conclusion} 化成了
与\eqref{eq:general-surface-index-formula} 相合的形式。
\end{proof}
到目前为止我们表述这些内容的方式有一个严重的缺点：$Gr(V)^\#$ 只是在相差%
\trHomotopyEquivalence{}的意义下定义的，这给抽象\trLinear\trLagrange
\trBrane{}的概念蒙上了一层模糊性。因此，我们将用一套\trEquivalence{}却更几何的框架——
\trGrading{}与\trPinStructure{}——来取代它。
